\documentclass[11pt]{article}
\usepackage[T1]{fontenc}
\usepackage[margin=1in]{geometry}
\usepackage{amsmath,amssymb,amsthm,mathtools,microtype,graphicx}
\usepackage[colorlinks=true,urlcolor=blue,citecolor=blue,linkcolor=blue]{hyperref}
\newtheorem{theorem}{Theorem}
\newtheorem{lemma}[theorem]{Lemma}
\newcommand{\R}{\mathbb R}
\newcommand{\B}{\mathbb B}
\newcommand{\vol}{\operatorname{vol}}
\newcommand{\Id}{\operatorname{Id}}
\newcommand{\Span}{\operatorname{span}}

\title{A family of counterexamples to the cube-constant conjecture for John ellipsoids of revolution}
\author{Counterexample packet for \texttt{fa\_banach\_001}}
\date{11 August 2026}

\begin{document}
\maketitle

\noindent\textbf{Status.}
Candidate full disproof of Conjecture~5.10 in arXiv:2507.06947, subject to
human review.  For every pair \(1\leq s<d\) for which \(s\nmid d\), an
explicit centrally symmetric polytope violates the proposed volume bound.
The smallest example is \((d,s)=(3,2)\).

\section*{Source conjecture}

G.~Ivanov, Z.~L\'angi, M.~Nasz\'odi, and \`A.~Sagmeister,
\emph{John Ellipsoids of Revolution}, arXiv:2507.06947v2 (5 August 2025)
\cite{source}.  Conjecture~5.10 on printed and PDF page~19 states that if
an origin-centered ellipsoid \(E\) solves the fixed-axis problem for a
centrally symmetric convex body \(K\subset\R^d\) and an \(s\)-dimensional
subspace \(F\), then
\begin{equation}\label{eq:conjecture}
 \frac{\vol_s(K\cap F)}{\vol_s(E\cap F)}
 \leq C_{\square}(d,s)
 \frac{\vol_s([-1,1]^s)}{\vol_s(\B^s)},
\end{equation}
where, on writing \(d=\ell s+t\), \(0\leq t<s\),
\begin{equation}\label{eq:cubeconstant}
 C_{\square}(d,s)^2=(\ell+1)^t\ell^{s-t}.
\end{equation}

\begin{center}
\includegraphics[width=0.98\linewidth]{figures/open_problem_crop.png}
\end{center}

\section*{Definitions and normalization}

Write \(\B^d=\{x\in\R^d:\lVert x\rVert_2\leq1\}\).  In the source,
an \(F\)-operator preserves both \(F\) and \(F^\perp\), acts as a nonzero
scalar on \(F^\perp\), and is arbitrary and invertible on \(F\).  Thus the
identity is an \(F\)-operator, and \(\B^d\) is an admissible ellipsoid of
revolution for every \(F\).  We prove below that it is largest even among
\emph{all} ellipsoids in the constructed body, which is stronger than the
hypothesis required in~\eqref{eq:conjecture}.

For vectors \(u,v\in\R^d\), \(u\otimes v\) denotes the rank-one operator
\(x\mapsto\langle x,v\rangle u\).

\section*{Proof intuition}

The conjectured constant contains an integer partition of \(d\): its square
is the largest product of \(s\) positive integer block sizes summing to
\(d\).  John's contact decomposition has positive \emph{real} weights and
does not impose this integrality.  We divide the total contact weight equally
among the \(s\) coordinate directions of \(F\), assigning the real weight
\(d/s\) to each.  Suitable lifts into \(F^\perp\) make the lifted unit normals
an exact John decomposition of the identity.  The resulting slab polytope
has section
\(\sqrt{d/s}\,[-1,1]^s\), so it attains the larger continuous
Brascamp--Lieb bound from the source.  When \(s\nmid d\), strict AM--GM says
that this continuous value exceeds the integer cube constant.

\section*{Counterexample theorem}

\begin{theorem}\label{thm:counterexample}
Let \(d\geq3\) and \(1\leq s<d\), and assume \(s\nmid d\).  There are an
\(s\)-dimensional subspace \(F\subset\R^d\) and a centrally symmetric
convex polytope \(K\subset\R^d\) such that:
\begin{enumerate}
\item \(\B^d\subset K\), and \(\B^d\) is a largest-volume ellipsoid in
      \(K\); in particular \(E=\B^d\) solves the source's fixed-axis
      ellipsoid-of-revolution problem for \(K\) and \(F\);
\item \(K\cap F=\sqrt{d/s}\,[-1,1]^s\) after identifying \(F\) with
      \(\R^s\);
\item
\[
 \frac{\vol_s(K\cap F)}{\vol_s(E\cap F)}
 > C_{\square}(d,s)
 \frac{\vol_s([-1,1]^s)}{\vol_s(\B^s)}.
\]
\end{enumerate}
Consequently Conjecture~5.10 is false for every nondivisible pair
\(1\leq s<d\).
\end{theorem}

\begin{proof}
Put \(r=d-s\), and choose orthonormal bases
\(\{e_1,\dots,e_s\}\) of \(F\) and
\(\{f_1,\dots,f_r\}\) of \(F^\perp\).  Set
\[
 p=\sqrt{\frac{s}{d}},\qquad q=\sqrt{\frac{r}{d}}.
\]
For \(1\leq j\leq s\), \(1\leq k\leq r\), and
\(\varepsilon,\delta\in\{-1,1\}\), define
\[
 u_{jk}^{\varepsilon,\delta}
 =\varepsilon p e_j+\delta q f_k.
\]
Every such vector is a unit vector.  Let
\begin{equation}\label{eq:K}
 K=\bigcap_{j,k,\varepsilon,\delta}
 \left\{x\in\R^d:
 \bigl|\langle x,u_{jk}^{\varepsilon,\delta}\rangle\bigr|\leq1
 \right\}.
\end{equation}
This is centrally symmetric.  The normals span \(\R^d\), since sums and
differences of the two choices of \(\delta\) yield multiples of \(e_j\)
and \(f_k\).  Hence \(K\) is bounded and is a convex body.  Since every
normal is unit, Cauchy--Schwarz gives \(\B^d\subset K\).

We next prove the maximal-ellipsoid assertion directly.  Give every indexed
normal the weight
\[
 \alpha=\frac{d}{4sr}.
\]
The sum of the weights is \(d\), and the sign symmetry makes the weighted
vector sum zero.  For fixed \(j,k\), cancellation of the cross terms gives
\[
 \sum_{\varepsilon,\delta}
 u_{jk}^{\varepsilon,\delta}\otimes
 u_{jk}^{\varepsilon,\delta}
 =4p^2 e_j\otimes e_j+4q^2 f_k\otimes f_k.
\]
Therefore
\begin{align*}
 \sum_{j,k,\varepsilon,\delta}
 \alpha\,u_{jk}^{\varepsilon,\delta}\otimes
 u_{jk}^{\varepsilon,\delta}
 &=\frac{d p^2}{s}\sum_{j=1}^s e_j\otimes e_j
   +\frac{d q^2}{r}\sum_{k=1}^r f_k\otimes f_k\\
 &=\Id_F+\Id_{F^\perp}=\Id_{\R^d}.
\end{align*}

For completeness, let \(z+T\B^d\subset K\) be any translated ellipsoid,
where \(T\) is invertible.  The two support inequalities for each symmetric
slab in~\eqref{eq:K} imply
\[
 |\langle z,u_{jk}^{\varepsilon,\delta}\rangle|
 +\lVert T^*u_{jk}^{\varepsilon,\delta}\rVert_2\leq1,
\]
and hence \(\lVert T^*u_{jk}^{\varepsilon,\delta}\rVert_2\leq1\).
Using the identity decomposition,
\[
 \operatorname{tr}(TT^*)
 =\sum_{j,k,\varepsilon,\delta}
   \alpha\lVert T^*u_{jk}^{\varepsilon,\delta}\rVert_2^2
 \leq\sum_{j,k,\varepsilon,\delta}\alpha=d.
\]
AM--GM applied to the squared singular values of \(T\) yields
\[
 |\det T|^{2/d}\leq\frac{\operatorname{tr}(TT^*)}{d}\leq1.
\]
Thus \(\vol_d(z+T\B^d)\leq\vol_d(\B^d)\), proving that \(\B^d\) is a
largest-volume ellipsoid in \(K\).

If \(x=\sum_{j=1}^s x_j e_j\in F\), every inequality in~\eqref{eq:K}
reduces to \(p|x_j|\leq1\).  Hence
\[
 K\cap F=\sqrt{\frac{d}{s}}[-1,1]^s,
\]
and consequently
\begin{equation}\label{eq:ratio}
 \frac{\vol_s(K\cap F)}{\vol_s(\B^d\cap F)}
 =\left(\frac{d}{s}\right)^{s/2}
  \frac{\vol_s([-1,1]^s)}{\vol_s(\B^s)}.
\end{equation}

Finally write \(d=\ell s+t\), where \(0<t<s\) because \(s\nmid d\).
Strict AM--GM applied to \(t\) copies of \(\ell+1\) and \(s-t\) copies of
\(\ell\) gives
\[
 \bigl((\ell+1)^t\ell^{s-t}\bigr)^{1/s}
 <\frac{t(\ell+1)+(s-t)\ell}{s}
 =\frac{d}{s}.
\]
Together with~\eqref{eq:cubeconstant}, this says
\(C_{\square}(d,s)<(d/s)^{s/2}\).  Comparing with~\eqref{eq:ratio}
proves the strict violation.
\end{proof}

\section*{Smallest counterexample}

Take \(d=3\), \(s=2\), \(F=\Span\{e_1,e_2\}\),
\(p=\sqrt{2/3}\), and \(q=1/\sqrt3\).  Removing duplicate symmetric slabs,
the body can be written
\[
 K=\left\{x\in\R^3:
 \begin{array}{ll}
 |p x_1+q x_3|\leq1,&|p x_1-q x_3|\leq1,\\
 |p x_2+q x_3|\leq1,&|p x_2-q x_3|\leq1
 \end{array}\right\}.
\]
The proof above shows that \(\B^3\) is its John ellipsoid, while
\[
 K\cap F=\sqrt{\frac32}[-1,1]^2,
 \qquad \operatorname{area}(K\cap F)=6.
\]
But \(C_{\square}(3,2)=\sqrt2\), so the conjectured unnormalized section
bound is \(4\sqrt2<6\).

\section*{Verification notes}

\begin{itemize}
\item The maximality proof includes arbitrary translations and arbitrary
      invertible \(T\); it does not merely verify the necessary contact
      equations or local maximality.
\item The construction gives a full John decomposition of \(\Id_{\R^d}\),
      stronger than the projected decomposition required for the constrained
      fixed-axis problem.
\item The supplied checker verifies unit norms, spanning, the weighted John
      identity, and the strict constant comparison for all \(d\leq15\).  It
      is a sanity check; the displayed identities are the proof.
\item The conjecture remains sharp on the divisible pairs within this family:
      if \(s\mid d\), then \(C_{\square}(d,s)=(d/s)^{s/2}\).
\end{itemize}

\section*{Novelty check and limitations}

A bounded search on 11 August 2026 checked the exact paper title, arXiv id,
``Conjecture 5.10'', ``John ellipsoids of revolution counterexample'', and
close volume-ratio phrases on the arXiv domain.  It recovered the source and
the cube-section paper that supplies \(C_{\square}\) \cite{cube}, but no later
resolution or counterexample.  The local run indexes likewise contained no
duplicate.  This supports candidate novelty only; specialist author/citation
review remains appropriate.

The theorem disproves the conjecture as written.  It does not determine the
sharp replacement constant for every centrally symmetric body, although the
construction shows that no universal constant smaller than
\((d/s)^{s/2}\) is possible.  Since the source already proves the matching
upper bound \((d/s)^{s/2}\), the same construction also identifies that bound
as sharp for every \(1\leq s<d\).

\section*{Human-review recommendation}

Validity confidence is high.  The main review points are: verify the four-sign
outer-product calculation; check the support-function inequality for a
translated ellipsoid; and confirm from the source definition that the unit
ball is an admissible \(F\)-ellipsoid.  Novelty confidence is moderate because
the counterexample is elementary and should be checked with the source
authors or a comprehensive citation search before publication.

\begin{thebibliography}{9}
\bibitem{source}
G.~Ivanov, Z.~L\'angi, M.~Nasz\'odi, and \`A.~Sagmeister,
\emph{John Ellipsoids of Revolution},
\href{https://arxiv.org/abs/2507.06947}{arXiv:2507.06947v2} (2025).

\bibitem{cube}
G.~Ivanov and I.~Tsiutsiurupa,
\emph{On the volume of sections of the cube},
Analysis and Geometry in Metric Spaces \textbf{9} (2021), 1--18;
\href{https://arxiv.org/abs/2004.02674}{arXiv:2004.02674}.
\end{thebibliography}

\end{document}
